Una massa de 0,5~kg penjada d'una molla segueix un moviment harmònic simple vertical. Una segona massa, subjectada a 19~cm del centre d'un disc, rota amb una velocitat angular de 6,41~rad/s. Quan les masses s'il·luminen lateralment, les dues ombres segueixen el mateix moviment.

\figura[0.45]{fig1.png}

\begin{apartat}{1,25}
Escriviu l'equació de la posició vertical de les ombres respecte al temps suposant que a $t = 0$ es troben en el punt més baix. Trobeu l'energia mecànica de la massa penjada de la molla i la constant elàstica d'aquesta.
\end{apartat}

\begin{apartat}{1,25}
Calculeu l'equació de la velocitat i de l'energia cinètica respecte al temps de la massa que penja de la molla. Representeu a la quadrícula adjunta l'energia mecànica, l'energia potencial i l'energia cinètica en funció de la posició vertical per a la massa que penja de la molla.
\end{apartat}

\figura[0.75]{fig2.png}
